\documentclass[12pt,spanish,mexico]{article}
\usepackage[utf8]{inputenc}
\PassOptionsToPackage{hyphens}{url}
\usepackage{hyperref}
\usepackage[es-nodecimaldot]{babel} % nombre de opción corregido
\usepackage{amsthm, amssymb, amsmath, amsfonts, bbm, mathtools}
\usepackage{csquotes}
\usepackage{enumerate}
\usepackage[shortlabels]{enumitem}
\usepackage{titling}
\usepackage{blindtext}
\usepackage{lmodern} 
\usepackage{float}
\usepackage{mathrsfs}
\usepackage[capitalise]{cleveref}
\usepackage{caption} 
\usepackage{graphicx}
\usepackage{chngcntr}
\usepackage{enumitem}    
\usepackage{cleveref} 
\usepackage[ruled,vlined]{algorithm2e}

\counterwithin{figure}{section}
\numberwithin{equation}{section}
\usepackage[vmargin=1in, hmargin=0.5in]{geometry}  

\setlength{\topmargin}{-50pt}
\setlength{\headheight}{38.23112pt}
\setlength{\headsep}{24pt}

\makeatletter
\def\iftitle{\ifx\thetitle \empty \else }
\def\safetitle{\iftitle{\thetitle}\fi}
\def\headertitle{%
    \ifdefined\hw{\hw}\fi%
    \ifdefined\hw \iftitle{ -- }\fi%
    \safetitle%
}
\def\@maketitle{%
  \begin{center}%
Tarea 5 - Problemas inversos - 15/04/2026 \\
Muñoz Macías Christian Jaime 
    \vskip 1em%
    \hrule
  \end{center}%
  \par
  \vskip 1.5em%
}
\makeatother
\usepackage{amsmath}
\usepackage{tikz}

\newcommand\RADot{%
  \mathord{%
    \mspace{1mu}%
    \text{\radot}%
    \mspace{1mu}%
  }%
}

\newcommand\radot{%
    \tikz[line cap=round,x=1ex,y=1ex,line width=0.3pt]
    {\draw (0,1) |- (1,0) (0.55,0) arc(0:90:0.55); \fill (0.23,0.23) circle (0.05);}%
}

%% encabezado personalizado
\usepackage{fancyhdr}
%% estilo plain personalizado
\fancypagestyle{plain}{%
    \fancyhf{} % limpia encabezado y pie de página
    \renewcommand{\headrulewidth}{0pt}
    \setlength{\headheight}{0pt}
    \fancyfoot[R]{\thepage}
}
%% estilo para todas las páginas
\pagestyle{fancy}
\fancyhf{}
\fancyhead[C]{Maestría en Matemáticas Aplicadas}
\fancyfoot[R]{\thepage}

%%%%%%%%%%%%%%%%%%%%%%%%
%% DEFINICIONES ÚTILES %%
%%%%%%%%%%%%%%%%%%%%%%%%
%% definiciones matemáticas
\newcommand{\N}{\mathbb{N}}
\newcommand{\Z}{\mathbb{Z}}
\newcommand{\Ha}{\mathcal{H}}
\newcommand{\R}{\mathbb{R}}
\newcommand{\Q}{\mathbb{Q}}
\newcommand{\C}{\mathbb{C}}
\newcommand{\PP}{\mathbb{P}}

\newtheorem{proposition}{Proposición}
\newtheorem{theorem}{Teorema}
\newtheorem{lemma}{Lema}
\newtheorem{corollary}{Corolario}

\theoremstyle{definition}
\newtheorem{problem}{Problema}
\newtheorem{definition}{Definición}[section]
\newtheorem*{example}{Ejemplo}

\theoremstyle{remark}
\newtheorem*{observation}{Observación}
\newtheorem*{note}{Nota}

\newenvironment{solution}{\begin{proof}[Solución]}{\end{proof}}
\newenvironment{dem}{\begin{proof}[Demostración]}{\end{proof}}
\renewcommand{\qedsymbol}{$\blacksquare$}

%% comandos útiles
\newcommand{\eps}{\varepsilon}
\newcommand{\To}{\longrightarrow}
\newcommand{\Eu}{\varPhi}
\newcommand{\Mapsto}{\longmapsto}

%% delimitadores emparejados
\DeclarePairedDelimiter{\set}{\{}{\}}
\DeclarePairedDelimiter{\floor}\lfloor\rfloor
\DeclarePairedDelimiter{\ceil}\lceil\rceil
\DeclarePairedDelimiter{\abs}\lvert\rvert
\DeclarePairedDelimiter{\gen}\langle\rangle
\DeclareMathOperator{\sgn}{sgn}
\DeclareMathOperator{\ord}{ord}

%% operador
\newcommand{\op}[1]{\operatorname{#1}}

%% etiquetas del documento
\title{Tarea 5 - Problemas inversos - 15/04/2026}
\def\course{Materia}
\def\university{CIMAT - MsC. Maestría en matemáticas aplicadas}
\author{Muñoz Macías Christian Jaime}
\def\email{\url{christian.munoz@cimat.mx}}
\date{\today}

%% permitir saltos de línea en ecuaciones
\allowdisplaybreaks[4]

\begin{document}

\maketitle


\section*{Modelo de una pandemia de influenza}

Dada una población de tamaño $N$, el contagio de influenza durante un brote pandémico está gobernado por el sistema de ecuaciones diferenciales

\[
\begin{align*}
\frac{dS}{dt} &= -\beta \frac{I}{N} S, \\
\frac{dE}{dt} &= \beta \frac{I}{N} S - \sigma E, \\
\frac{dI}{dt} &= \sigma E - \gamma I, \\
\frac{dR}{dt} &= \gamma I,
\end{align*}
\]

con condiciones iniciales

\[
\begin{align*}
S(0) &= S_0 \\
E(0) &= E_0 \\
I(0) &= I_0 \\
R(0) &= R_0
\end{align*}
\]

Las variables $S(t)$, $E(t)$, $I(t)$ y $R(t)$ denotan el número de personas susceptibles, latentes, infecciosas, y removidas respectivamente al tiempo $t$. Aquí, removidos significa que las personas no participan en el contagio por cualquier razón: aislamiento, muerte, recuperación, etc. El número diario de nuevos casos $c_i$ en el día $i$ está dado por

\[
c_i = q \int_{t_{i-1}}^{t_i} \sigma E(s) \, ds, \tag{3}
\]

donde $[t_{i-1}, t_i]$ es el plazo entre dos tiempos de observación, que en el presente caso es igual a un día. Los parámetros y condiciones iniciales del modelo (1)-(2) son como se indica en la siguiente tabla:

\begin{center}
\begin{tabular}{|l|l|c|}
\hline
\textbf{Nombre} & \textbf{Significado} & \textbf{Valor} \\
\hline
$\sigma$ & Tasa de latencia & $1/3$ días$^{-1}$ \\
\hline
$\gamma$ & Tasa de recuperación & $1/5$ días$^{-1}$ \\
\hline
$q$ & probabilidad de observación de una infección & $0.5$ \\
\hline
$\beta$ & Tasa de contacto & por inferir \\
\hline
$S_0$ & Número inicial de susceptibles & por inferir \\
\hline
$E_0$ & Número inicial de latentes & por inferir \\
\hline
$I_0$ & Número inicial de infecciosos & por inferir \\
\hline
$R_0$ & Número inicial de removidos & por inferir \\
\hline
\end{tabular}
\end{center}

Resuelve un problema de asimilación de datos tomando en cuenta lo siguiente:

\begin{enumerate}
\item Suponemos $N = 550,000$.
\item La condición $N = S_0 + E_0 + I_0 + R_0$ reduce en uno el número de parámetros por estimar.
\item Consideramos el archivo de datos \texttt{pandemic\_influenza\_SF\_1918.xls} que se encuentra en la carpeta de Dropbox. Estos datos corresponden a la primera ola de la pandemia de influenza de 1918 en San Francisco, California, Estados Unidos.
\item Usa el método de tu elección para estimar los parámetros del sistema dinámico (1)-(2) con el modelo de observación (3).
\item Interpreta las condiciones iniciales en términos de cuanta gente se aisló al inicio de la pandemia.
\end{enumerate}

\begin{solution}

El sistema dinámico se escribe como
\[
\frac{dy}{dt}=f(y),
\qquad
y=
\begin{pmatrix}
S\\E\\I\\R
\end{pmatrix},
\]

con
\[
f(y)=
\begin{pmatrix}
-\beta \frac{SI}{N}\\
\beta \frac{SI}{N}-\sigma E\\
\sigma E-\gamma I\\
\gamma I
\end{pmatrix}.
\]

Por el teorema fundamental del cálculo,
\[
y(t+h)-y(t)=\int_t^{t+h} f(y(\tau))\,d\tau.
\]


Se realiza expansión de Taylor en $y(\tau)$ alrededor de $y(t)$:

\[
f(y(\tau))
=
f(y(t)) + J(y(t))\big(y(\tau)-y(t)\big) + O(\|y(\tau)-y(t)\|^2),
\]

donde el Jacobiano es
\[
J(y)=\frac{\partial f}{\partial y}(y)\in\mathbb{R}^{4\times 4}.
\]

Explícitamente:
\[
J(y)=
\begin{pmatrix}
-\beta \frac{I}{N} & 0 & -\beta \frac{S}{N} & 0 \\
\beta \frac{I}{N} & -\sigma & \beta \frac{S}{N} & 0 \\
0 & \sigma & -\gamma & 0 \\
0 & 0 & \gamma & 0
\end{pmatrix}.
\]

Se usa expansión de Taylor en el tiempo:
\[
y(\tau)=y(t)+(\tau-t)y'(t)+O((\tau-t)^2),
\]

y como $y'(t)=f(y(t))$, entonces
\[
y(\tau)-y(t)=(\tau-t)f(y(t))+O(h^2).
\]


Sustituyendo:

\[
f(y(\tau))
=
f(y(t)) + (\tau-t)J(y(t))f(y(t)) + O(h^2).
\]

Entonces
\[
\int_t^{t+h} f(y(\tau))\,d\tau
=
\int_t^{t+h} f(y(t))\,d\tau
+
\int_t^{t+h} (\tau-t)J(y(t))f(y(t))\,d\tau
+
O(h^3).
\]


Primera integral:
\[
\int_t^{t+h} f(y(t))\,d\tau = h f(y(t)).
\]

Segunda integral:
\[
\int_t^{t+h} (\tau-t)\,d\tau = \frac{h^2}{2}.
\]

Por tanto:
\[
\int_t^{t+h} f(y(\tau))\,d\tau
=
h f(y(t))
+
\frac{h^2}{2}J(y(t))f(y(t))
+
O(h^3).
\]


\[
y(t+h)-y(t)
=
h f(y(t))
+
\frac{h^2}{2}J(y(t))f(y(t))
+
O(h^3).
\]

Equivalente en forma discreta:
\[
y_{n+1}
=
y_n
+
h f(y_n)
+
\frac{h^2}{2}J(y_n)f(y_n).
\]

El número de casos observados por día se modela como:
\[
c_i = q \int_{t_{i-1}}^{t_i} \sigma E(s)\,ds.
\]

Definimos el modelo predictivo:
\[
c_i^{mod}(\theta) = q \int_{t_{i-1}}^{t_i} \sigma E_{mod}(s;\theta)\,ds,
\]
donde
\[
\theta = (\beta, S_0, E_0, I_0, R_0).
\]

Se conoce que
\[
N = S_0 + E_0 + I_0 + R_0,
\]
entonces:
\[
S_0 = N - (E_0 + I_0 + R_0).
\]

Por lo tanto, los parámetros a estimar son:
\[
\theta = (\beta, E_0, I_0, R_0).
\]

Se introduce la reparametrización:
\[
\theta = (\beta, u_1, u_2, u_3), \qquad u_i > 0,
\]
y se define
\[
D = 1 + u_1 + u_2 + u_3.
\]

Entonces:
\[
S_0 = \frac{N}{D}, \quad
E_0 = \frac{N u_1}{D}, \quad
I_0 = \frac{N u_2}{D}, \quad
R_0 = \frac{N u_3}{D}.
\]

Con esta parametrización se satisface:
\[
S_0 + E_0 + I_0 + R_0 = N, \qquad S_0,E_0,I_0,R_0 \ge 0.
\]

Por lo tanto:
\[
c_i^{mod}(\theta) = c_i^{mod}(\beta, u_1, u_2, u_3).
\]

Se asume ruido gaussiano:
\[
c_i = c_i^{mod}(\theta) + \varepsilon_i,
\qquad
\varepsilon_i \sim \mathcal{N}(0,\eta).
\]

La verosimilitud es:
\[
p(c \mid \theta)
\propto
\exp\left(
-\frac{1}{2\eta}\sum_{i=1}^T (c_i - c_i^{mod}(\theta))^2
\right).
\]

Su versión logarítmica:
\[
\log p(c \mid \theta)
=
-\frac{1}{2\eta}\sum_{i=1}^T (c_i - c_i^{mod}(\theta))^2 + C.
\]

Para la priori se usan variables positivas mediante transformación logarítmica:
\[
\tilde{\theta} = (\log \beta, \log u_1, \log u_2, \log u_3).
\]

Se asume:
\[
\tilde{\theta} \sim \mathcal{N}(\mu, \Sigma),
\]
es decir:
\[
p(\tilde{\theta})
\propto
\exp\left(
-\frac{1}{2}(\tilde{\theta}-\mu)^T \Sigma^{-1} (\tilde{\theta}-\mu)
\right).
\]

Por cambio de variable:
\[
p(\theta)
=
p(\tilde{\theta})
\left|
\det\left(\frac{\partial \tilde{\theta}}{\partial \theta}\right)
\right|.
\]

Como \(\tilde{\theta}_i = \log \theta_i\), el Jacobiano es:
\[
J =
\begin{pmatrix}
\frac{1}{\beta} & 0 & 0 & 0 \\
0 & \frac{1}{u_1} & 0 & 0 \\
0 & 0 & \frac{1}{u_2} & 0 \\
0 & 0 & 0 & \frac{1}{u_3}
\end{pmatrix},
\qquad
|\det(J)| = \frac{1}{\beta u_1 u_2 u_3}.
\]

Entonces:
\[
\log p(\theta)
=
\log p(\tilde{\theta})
-
\sum_{i=1}^4 \tilde{\theta}_i.
\]

Finalmente:
\[
\log p(\theta)
=
-\frac{1}{2}(\tilde{\theta}-\mu)^T \Sigma^{-1} (\tilde{\theta}-\mu)
-\sum_{i=1}^4 \tilde{\theta}_i + C.
\]

La posterior queda:
\[
p(\theta \mid c) \propto p(c \mid \theta)\, p(\theta),
\]

y su versión logarítmica:
\[
\log p(\theta \mid c)
=
-\frac{1}{2\eta}\sum_{i=1}^T (c_i - c_i^{mod}(\theta))^2
-\frac{1}{2}(\tilde{\theta}-\mu)^T \Sigma^{-1} (\tilde{\theta}-\mu)
-\sum_{i=1}^4 \tilde{\theta}_i + C.
\]
El parámetro \(\eta > 0\), el cual representa la varianza del error en las observaciones y del desajuste del modelo entre los datos y la solución numérica del sistema. Se asume ruido gaussiano:
\[
\varepsilon_i \sim \mathcal{N}(0,\eta).
\]

Por lo tanto, el modelo de observación se escribe como:
\[
c_i = c_i^{mod}(\theta) + \varepsilon_i.
\]

El parámetro \(\eta\) controla la magnitud del error cuadrático medio entre datos y modelo:
\[
\eta \approx \mathbb{V}ar(c_i).
\]

La matriz de covarianza \(\Sigma\) del prior en el espacio transformado se propone como:
\[
\Sigma =0.01 \mathrm{diag}\left(1,1,1,1\right).
\]
\section{Resultados}

Los parámetros estimados mediante MCMC son:

\[
\begin{aligned}
\beta &= 5.4003 \pm 0.0423 \\
u_E &= 0.0126 \pm 0.0052 \\
u_I &= 0.4324 \pm 0.0943 \\
u_R &= 0.0104 \pm 0.0070
\end{aligned}
\]

donde las incertidumbres corresponden a dos desviaciones estándar (\( \pm 2\sigma \)).

\subsection{Evolución de las cadenas MCMC}

A continuación se muestra la trayectoria de la cadena para el parámetro de transmisión \( \beta \). Las líneas discontinuas y punteadas indican la media y la banda de \( \pm 2\sigma \), respectivamente.

\begin{figure}[H]
    \centering
    \includegraphics[width=0.6\linewidth]{C:/Users/chris/OneDrive/Escritorio/python/CIMAT/problemas inversos/Tarea_5/Cadena MCMC - beta.pdf}
    \caption{Evolución de la cadena MCMC para \( \beta \). La línea horizontal representa la media posterior.}
\end{figure}

Las siguientes figuras presentan las cadenas correspondientes a los parámetros de condiciones iniciales \( u_E \), \( u_I \) y \( u_R \).

\begin{figure}[H]
    \centering
    \includegraphics[width=0.6\linewidth]{C:/Users/chris/OneDrive/Escritorio/python/CIMAT/problemas inversos/Tarea_5/Cadena MCMC - u_E.pdf}
    \caption{Evolución de la cadena MCMC para \( u_E \).}
\end{figure}

\begin{figure}[H]
    \centering
    \includegraphics[width=0.6\linewidth]{C:/Users/chris/OneDrive/Escritorio/python/CIMAT/problemas inversos/Tarea_5/Cadena MCMC - u_I.pdf}
    \caption{Evolución de la cadena MCMC para \( u_I \).}
\end{figure}

\begin{figure}[H]
    \centering
    \includegraphics[width=0.6\linewidth]{C:/Users/chris/OneDrive/Escritorio/python/CIMAT/problemas inversos/Tarea_5/Cadena MCMC - u_R.pdf}
    \caption{Evolución de la cadena MCMC para \( u_R \).}
\end{figure}

\subsection{Distribuciones posteriores}

La figura muestra las distribuciones marginales y las correlaciones entre pares de parámetros. Los histogramas en la diagonal representan las densidades posteriores marginales, mientras que los gráficos de dispersión fuera de la diagonal ilustran las correlaciones posteriores.

\begin{figure}[H]
    \centering
    \includegraphics[width=0.6\linewidth]{C:/Users/chris/OneDrive/Escritorio/python/CIMAT/problemas inversos/Tarea_5/distribucion_pa.pdf}
    \caption{Matriz de distribuciones posteriores. Diagonal: densidades marginales. Fuera de diagonal: correlaciones entre parámetros.}
\end{figure}

\subsection{Dinámica de los compartimentos epidemiológicos}

Las siguientes figuras presentan la evolución temporal de los compartimentos del modelo SEIR: susceptibles, expuestos, infectados y recuperados. Las bandas sombreadas representan el intervalo de incertidumbre del 95\% (\( \pm 2\sigma \)) obtenido a partir de las muestras posteriores.

\begin{figure}[H]
    \centering
    \includegraphics[width=0.6\linewidth]{C:/Users/chris/OneDrive/Escritorio/python/CIMAT/problemas inversos/Tarea_5/Susceptibles.pdf}
    \caption{Evolución de la población susceptible \( S(t) \) con banda de incertidumbre del 95\%.}
\end{figure}

\begin{figure}[H]
    \centering
    \includegraphics[width=0.6\linewidth]{C:/Users/chris/OneDrive/Escritorio/python/CIMAT/problemas inversos/Tarea_5/Expuestos.pdf}
    \caption{Evolución de la población expuesta \( E(t) \) con banda de incertidumbre del 95\%.}
\end{figure}

\begin{figure}[H]
    \centering
    \includegraphics[width=0.6\linewidth]{C:/Users/chris/OneDrive/Escritorio/python/CIMAT/problemas inversos/Tarea_5/Infectados.pdf}
    \caption{Evolución de la población infectada \( I(t) \) con banda de incertidumbre del 95\%.}
\end{figure}

\begin{figure}[H]
    \centering
    \includegraphics[width=0.6\linewidth]{C:/Users/chris/OneDrive/Escritorio/python/CIMAT/problemas inversos/Tarea_5/Recuperados.pdf}
    \caption{Evolución de la población recuperada \( R(t) \) con banda de incertidumbre del 95\%.}
\end{figure}

\subsection{Incidencia modelada vs. datos observados}

La siguiente figura compara la incidencia reportada (puntos rojos) con la predicción del modelo SEIR (línea azul). La banda sombreada representa la incertidumbre del 95\% proveniente de la propagación de incertidumbre de los parámetros.

\begin{figure}[H]
    \centering
    \includegraphics[width=0.6\linewidth]{C:/Users/chris/OneDrive/Escritorio/python/CIMAT/problemas inversos/Tarea_5/Incidencia - Modelo SEIR con incertidumbre.pdf}
    \caption{Incidencia semanal: comparación entre datos observados y predicción del modelo. La banda sombreada indica el intervalo de incertidumbre del 95\%.}
\end{figure}

\subsection{Discusión}

Los resultados presentados exhiben una sensibilidad considerable a diversos aspectos del procedimiento de estimación:

\begin{enumerate}
    \item \textbf{Configuración del muestreador MCMC}: el tamaño de paso y la longitud de la cadena afectan la convergencia y la mezcla de las muestras.
    \item \textbf{Condiciones iniciales}: el punto de partida de los parámetros influye en la región del espacio paramétrico explorada.
    \item \textbf{Distribución a priori}: la elección de la media \( \mu \) y la matriz de covarianza \( \Sigma \) de la distribución prior tiene un impacto significativo en las estimaciones finales.
    \item \textbf{Método de optimización}: el uso de diferentes algoritmos de optimización para obtener el punto inicial puede conducir a resultados distintos.
\end{enumerate}

Esta sensibilidad sugiere que el problema de estimación es inherentemente no lineal y posiblemente mal condicionado. Las principales fuentes de incertidumbre son:

\begin{itemize}
    \item \textbf{Falta de información identificable} en los datos observados para determinar de manera única todos los parámetros del modelo.
    \item \textbf{Especificación de la verosimilitud}: la elección de \( \eta \) (varianza del error) debe revisarse, ya que un valor fijo puede no capturar adecuadamente la heterocedasticidad de los datos.
    \item \textbf{Matriz de covarianza de la priori}: \( \Sigma = 0.01I \) podría ser demasiado restrictiva o mal escalada para los parámetros transformados logarítmicamente.
\end{itemize}


\end{solution}
\end{document}